\documentclass[12pt]{article}
\usepackage[utf8]{inputenc}
\usepackage[T1]{fontenc}
\usepackage{hyperref}
\title{Dhuwal Language Resources for K-12 Curriculum Development}
\author{Compiled by Claude Sonnet 5.5 (claude-sonnet-5-5) for Andr\'as Kornai}
\date{}
\begin{document}
\maketitle

\section{Introduction}
Dhuwal (ISO 639-3: dwu) is an Australian Aboriginal language of north-east Arnhem Land, part of the wider Yolngu Matha continuum. Web searches turned up little material specific to this code; most available resources cover Yolngu Matha varieties generally, or neighbouring varieties such as Djambarrpuyngu and Gupapuyngu. Nothing found is a ready-made K-12 curriculum for Dhuwal.

\section{Primary Educational Resources}
No Dhuwal-specific primers, textbooks or school programmes were found in the searches. The Charles Darwin University Yolngu Studies page lists a Yolngu Matha app (sounds, grammar, kinship and stories, written in Gupapuyngu, a related variety): \url{https://yolngustudies.cdu.edu.au/gupaappdownload.html}.

\section{Linguistic and Reference Materials}
\begin{itemize}
\item Wikipedia overview of the Dhuwal language: \url{https://en.wikipedia.org/wiki/Dhuwal_language}
\item Jeffrey Heath, \emph{Dhuwal (Arnhem Land) texts on kinship and other subjects, with grammatical sketch and dictionary} (1980), hosted on Deep Blue: \url{https://deepblue.lib.umich.edu/items/a6721d13-a93d-471c-b05d-bf641a92f3bc}
\item English--Yolngu Matha Legal Dictionary (ARDS; translated into Djambarrpuyngu, a related variety; the publisher's copyright notice restricts copying): \url{https://ards.com.au/s/English-Yolngu-Matha-Legal-Dictionary.pdf}
\end{itemize}

\section{Religious and Community Resources}
None specific to Dhuwal were identified in the searches.

\section{Digital and Audio Resources}
The NT Health library lists a Yolnu Matha health resource (reported to include audio in Dhuwal, Dhuwala and Dhangu): \url{https://library.health.nt.gov.au/mobileresources/YolnuMatha}

\section{Curriculum Development Context}
The reference description says Dhuwal is not known to be taught in schools. Related Yolngu Matha varieties have bilingual programmes (see the Dhuwaya dossier), which may offer models. Any curriculum work should be done with Yolngu communities and educators; the evidence found here is thin.
\end{document}